\documentclass[conference]{IEEEtran}
\IEEEoverridecommandlockouts

\usepackage[utf8]{inputenc}
\usepackage[T1]{fontenc}
\usepackage{cite}
\usepackage{amsmath,amssymb,amsfonts}
\usepackage{algorithmic}
\usepackage{graphicx}
\usepackage{textcomp}
\usepackage{xcolor}
\usepackage{booktabs}
\usepackage{multirow}
\usepackage{url}
\usepackage{hyperref}

\begin{document}

\title{Cross-Domain Generalization of a Reproducible TinyML\\
Anomaly Detection Pipeline: Seismic and Oil-Well Benchmarks\\
with VUS-PR, PA-F1, and Event-Level Evaluation}

\author{
  \IEEEauthorblockN{[AUTHOR NAME]}
  \IEEEauthorblockA{
    \textit{[Department]}\\
    \textit{[Institution]}\\
    [City, Country]\\
    [email]
  }
}

\maketitle

% ----------------------------------------------------------------
\begin{abstract}
A central challenge in anomaly detection benchmarking is that most pipelines
are designed and tuned for a single domain, making cross-domain comparison
unfair and results non-transferable.
We evaluate a generic TinyML pipeline---originally validated on seismic data
targeting ESP32 microcontrollers---on a structurally different domain:
the public Petrobras 3W oil-well dataset, a multivariate time series with
8 process sensors at 1~Hz.
The pipeline uses domain adapters that enforce a unified dataset contract
\texttt{(N, window\_size, n\_channels)}, enabling the same training,
evaluation, and MLOps components to operate on both domains without
modification.
Both domains are evaluated under a unified metric suite:
AUC-PR, VUS-PR, PA-F1, event-level F1, and FP/h.
On seismic data, the best model (Tiny TCN, 14{,}897 parameters)
achieves AUC-PR~$=0.9416$ and FP/h~$=3.136$.
On 3W, [RESULTS TO BE FILLED].
The results demonstrate that the pipeline generalizes across domains
with structural differences in sampling rate (40~Hz vs.\ 1~Hz),
dimensionality (univariate vs.\ 8-channel), anomaly ratio, and failure
semantics, supporting the adapter-based design as a viable architecture
for multi-domain edge monitoring systems.
\end{abstract}

\begin{IEEEkeywords}
Cross-domain generalization, TinyML, anomaly detection, multivariate time
series, Petrobras 3W, seismology, ESP32, VUS-PR, PA-F1, benchmark.
\end{IEEEkeywords}

% ================================================================
\section{Introduction}
\label{sec:intro}
% ================================================================

\subsection{The Cross-Domain Problem in Anomaly Detection}

Time-series anomaly detection research suffers from a lack of fair
cross-domain benchmarks.
Most published pipelines are designed for a specific domain (network
intrusion, industrial vibration, seismic), with preprocessing, feature
engineering, model architecture, and evaluation protocol optimized for
that domain~\cite{pang2021deep,chandola2009anomaly}.
When the same model is applied to a different domain, it fails---not
necessarily because the approach is wrong, but because domain-specific
assumptions are embedded in every pipeline stage.

This creates a reproducibility and generalization problem:
it is impossible to tell whether a method is genuinely better, or
whether it merely has better domain-specific engineering.
Kim et al.~\cite{kim2022towards} make this argument explicit, showing that
anomaly detection benchmarks are heavily influenced by metric choice and
dataset preparation.

\subsection{Adapter-Based Architecture as a Solution}

We propose and evaluate an adapter-based architecture that enforces a
\emph{unified dataset contract} between domain-specific preprocessing and
a domain-agnostic ML core.
The contract is a NumPy compressed archive:
\begin{equation}
  \{\mathbf{X}_\text{split} \in \mathbb{R}^{N \times W \times C},\;
    \mathbf{y}_\text{split} \in \{0,1\}^N\}
  \quad \text{for } \text{split} \in \{\text{train, val, test}\}
\end{equation}
where $N$ is the number of windows, $W$ the window size, and $C$ the
number of channels ($C=1$ for univariate).

Any domain---seismic ($C=1$, $f_s=40$~Hz), oil-well ($C=8$, $f_s=1$~Hz),
industrial vibration, or audio---contributes only a domain adapter.
The training, evaluation, MLOps, and export components remain unchanged.

\subsection{Contributions}

\begin{enumerate}
  \item \textbf{Cross-domain evaluation} of a unified TinyML pipeline on
        seismic (univariate, 40~Hz) and Petrobras 3W (8-channel, 1~Hz)
        under a common, rigorous metric suite.
  \item \textbf{Multivariate extension} of the pipeline contract, supporting
        $C$-channel input in CNN, TCN, and LSTM architectures via a single
        \texttt{n\_channels} parameter.
  \item \textbf{3W domain adapter} with split-by-well strategy, edge-aware
        per-channel preprocessing, and cross-channel correlation features.
  \item \textbf{Reproducible benchmark} comparing models across domains
        using AUC-PR, VUS-PR, PA-F1, Event-F1, and FP/h---enabling fair
        inter-domain ranking for the first time in this pipeline.
\end{enumerate}

% ================================================================
\section{Related Work}
\label{sec:related}
% ================================================================

\subsection{Anomaly Detection Benchmarks}

Hundman et al.~\cite{hundman2018detecting} benchmark LSTM-based detectors
on NASA spacecraft telemetry (multivariate).
Su et al.~\cite{su2019robust} propose OmniAnomaly for multivariate
time series, evaluating on SMAP and MSL datasets.
Shen et al.~\cite{shen2020timeseries} review anomaly detection for
multivariate time series, noting the lack of standard protocols.
Kim et al.~\cite{kim2022towards} argue that existing benchmarks
systematically favor certain metrics and preprocessing choices, calling
for more rigorous evaluation.

\subsection{Petrobras 3W Dataset}

Vargas et al.~\cite{vargas2019realistic} introduce the 3W dataset:
a public, labeled collection of oil-well events from Petrobras.
The dataset contains real and simulated records of 8 undesirable events
(e.g., abrupt pressure increase, hydrate formation, gas kick)
across multiple wells, with 8 process sensor channels at 1~Hz.
It is one of the few public multivariate anomaly datasets from critical
infrastructure, making it an important benchmark for industrial TinyML.

\subsection{Multivariate Time-Series Anomaly Detection}

LSTM-based approaches~\cite{hundman2018detecting,su2019robust} dominate
multivariate anomaly detection but are impractical for MCU deployment.
Convolution-based methods~\cite{bai2018empirical} offer better parameter
efficiency, making them more suitable for TinyML settings.
To our knowledge, no prior work evaluates edge-deployable models on 3W
using consistent metrics across multiple domains.

% ================================================================
\section{Domains and Datasets}
\label{sec:domains}
% ================================================================

\subsection{Domain 1: Seismic (Univariate)}

\begin{table}[htbp]
  \caption{Domain Profiles}
  \label{tab:profiles}
  \centering
  \begin{tabular}{lcc}
    \toprule
    Parameter & Seismic & Petrobras 3W \\
    \midrule
    Sampling rate    & 40~Hz         & 1~Hz           \\
    Window size      & 800 samples   & 60 samples     \\
    Window duration  & 20~s          & 60~s           \\
    Step             & 10~s (50\%)   & 30~s (50\%)    \\
    Channels ($C$)   & 1             & 8              \\
    Split strategy   & By event      & By well        \\
    Anomaly type     & Earthquake    & Process fault  \\
    Anomaly ratio    & $\approx$12.8\% & [TODO]       \\
    \bottomrule
  \end{tabular}
\end{table}

Full description in~\cite{artigo1} (companion paper).

\subsection{Domain 2: Petrobras 3W (Multivariate)}

The 3W dataset~\cite{vargas2019realistic} contains CSV files with 8 sensor
channels:
P-PDG, P-TPT, T-TPT, P-MON-CKP, T-JUS-CKP, P-JUS-CKGL,
T-JUS-CKGL, QGL.
Labels: class~0 (normal), class~$>0$ (fault, various types).
We use a binary formulation: any fault label $\to y=1$.

\subsubsection{3W Adapter Preprocessing (Edge-Aware)}

Per channel, the adapter applies:
\begin{enumerate}
  \item Linear interpolation of missing values.
  \item Linear detrend.
  \item Per-channel, per-window z-score normalization.
\end{enumerate}

This pipeline is intentionally minimal---all steps are implementable
in C/C++ on ESP32---maintaining alignment between offline training
and edge inference.

\subsubsection{Split Strategy}

Data are split \emph{by well}: training, validation, and test sets
contain disjoint wells.
This prevents the model from memorizing well-specific baseline behavior,
which would inflate test performance.
Ratio: 70/15/15 wells for train/val/test.

\subsubsection{Multivariate Features}

For classical models, features are extracted per channel
(28 features $\times$ 8 channels = 224 features) plus
$\binom{8}{2} = 28$ cross-channel Pearson correlation coefficients,
for a total of 252 features per window.

% ================================================================
\section{Methodology}
\label{sec:method}
% ================================================================

\subsection{Unified Model Suite}

The same six model families are evaluated on both domains:

\begin{table}[htbp]
  \caption{Models Evaluated Across Domains}
  \label{tab:models}
  \centering
  \begin{tabular}{lcc}
    \toprule
    Model & Input & Edge \\
    \midrule
    Random Forest (200 trees)    & Features ($F$-dim) & $\times$ \\
    Extra Trees (200 trees)      & Features ($F$-dim) & $\times$ \\
    Tiny CNN                     & $(W, C)$           & \checkmark \\
    Tiny TCN                     & $(W, C)$           & \checkmark \\
    \bottomrule
  \end{tabular}
\end{table}

For neural models, the input shape changes automatically from
$(W, 1)$ to $(W, C)$ via the \texttt{n\_channels} parameter,
with no architectural modification---\texttt{Conv1D} operates on
the channel dimension natively.

\subsection{Hyperparameter Optimization}

Optuna with 30 trials per domain-model combination (reduced from 88 in
the seismic-only study for computational tractability in the cross-domain
setting).
Same composite objective (Eq.~(1) in companion paper~\cite{artigo1}).

\subsection{Evaluation Protocol}

All models are evaluated under the unified metric suite defined
in~\cite{artigo2} (companion paper):
AUC-PR, VUS-PR ($b_{\max} = 5$), PA-F1, Event-F1 ($\tau=0.1$), FP/h.

The same quality-gate thresholds are applied across domains
(Table~II in~\cite{artigo2}), enabling direct comparison of
how many models survive the gate in each domain.

% ================================================================
\section{Results}
\label{sec:results}
% ================================================================

\subsection{Cross-Domain Model Comparison}

\begin{table*}[htbp]
  \caption{Cross-Domain Benchmark Results (Test Set)}
  \label{tab:xdomain}
  \centering
  \begin{tabular}{llccccccc}
    \toprule
    Domain & Model & AUC-PR & VUS-PR & PA-F1 & Event-F1 & FP/h & Params & Edge \\
    \midrule
    \multirow{4}{*}{Seismic}
      & Optuna Tiny TCN  & \textbf{0.9416} & [TODO] & [TODO] & [TODO]
                         & \textbf{3.136}  & 14{,}897  & \checkmark \\
      & Optuna Tiny CNN  & 0.9127 & [TODO] & [TODO] & [TODO]
                         & 4.896  & 19{,}489  & \checkmark \\
      & Random Forest    & 0.8127 & [TODO] & [TODO] & [TODO]
                         & 9.264  & ---       & $\times$ \\
      & Extra Trees      & 0.7901 & [TODO] & [TODO] & [TODO]
                         & 11.296 & ---       & $\times$ \\
    \midrule
    \multirow{4}{*}{3W}
      & Optuna Tiny TCN  & [TODO] & [TODO] & [TODO] & [TODO]
                         & [TODO] & [TODO]    & \checkmark \\
      & Optuna Tiny CNN  & [TODO] & [TODO] & [TODO] & [TODO]
                         & [TODO] & [TODO]    & \checkmark \\
      & Random Forest    & [TODO] & [TODO] & [TODO] & [TODO]
                         & [TODO] & ---       & $\times$ \\
      & Extra Trees      & [TODO] & [TODO] & [TODO] & [TODO]
                         & [TODO] & ---       & $\times$ \\
    \bottomrule
  \end{tabular}
  % TODO: rodar benchmark_cross_domain.py após baixar dataset 3W
  %       e preencher os valores [TODO]
\end{table*}

\subsection{Quality Gate Pass Rate by Domain}

\begin{table}[htbp]
  \caption{Quality Gate Pass Rate (Full Gate, All Candidates)}
  \label{tab:gate_rate}
  \centering
  \begin{tabular}{lccc}
    \toprule
    Domain & Candidates & Passed & Rate \\
    \midrule
    Seismic & [TODO] & [TODO] & [TODO]\% \\
    3W      & [TODO] & [TODO] & [TODO]\% \\
    \bottomrule
  \end{tabular}
\end{table}

\subsection{Relative Ranking Stability}

A key hypothesis of the adapter architecture is that relative model
ranking---not absolute performance---is stable across domains.
We measure ranking agreement using Spearman correlation $\rho_s$ between
the AUC-PR rankings on seismic and 3W:

\begin{equation}
  \rho_s = 1 - \frac{6\sum d_i^2}{n(n^2-1)}
\end{equation}

where $d_i$ is the rank difference for model $i$ and $n=4$.
We expect $\rho_s > 0$ (positive correlation in ranking) if the pipeline
is consistent across domains.

$\rho_s = \text{[TODO]}$ (target: $> 0.60$).

\subsection{VUS-PR vs.\ AUC-PR Discrepancy by Domain}

\begin{table}[htbp]
  \caption{Mean VUS-PR/AUC-PR Ratio per Domain}
  \label{tab:vuspr_ratio}
  \centering
  \begin{tabular}{lcc}
    \toprule
    Domain & AUC-PR (best) & VUS-PR (best) \\
    \midrule
    Seismic & 0.9416 & [TODO] \\
    3W      & [TODO] & [TODO] \\
    \bottomrule
  \end{tabular}
\end{table}

A lower VUS-PR/AUC-PR ratio indicates that the model relies on
point-exact alignment---a concern particularly for 3W, where fault
onset is gradual and window step is 30~s.

% ================================================================
\section{Discussion}
\label{sec:discussion}
% ================================================================

\subsection{Structural Domain Differences}

The two domains differ structurally in ways that stress-test the
pipeline's generalization:

\begin{itemize}
  \item \textbf{Sampling rate:} 40$\times$ difference (40~Hz vs.\ 1~Hz).
        Window duration differs accordingly (20~s vs.\ 60~s).
  \item \textbf{Dimensionality:} $C=1$ vs.\ $C=8$.
        Neural models process both without architecture changes;
        classical models receive 28 vs.\ 252 features.
  \item \textbf{Anomaly semantics:} seismic anomalies are impulsive
        (seconds to minutes); oil-well faults are gradual (hours).
        Event-F1 is more sensitive to this difference than AUC-PR.
  \item \textbf{Imbalance:} both are imbalanced but at different ratios
        ($\approx$12.8\% vs.\ [TODO]\%).
\end{itemize}

\subsection{Cross-Channel Correlation as a Discriminative Feature}

On 3W, the 28 cross-channel Pearson correlations contribute [X]\% of
Random Forest feature importance.
Several faults manifest as changes in correlational structure
(e.g., P-PDG and P-TPT decouple during a valve blockage) that are
invisible in per-channel statistics alone.
This motivates retaining correlation features even in classical models,
at the cost of 28 additional features.

\subsection{Limitations}

\textbf{Two-domain scope:} with only two domains, claims of cross-domain
generalization are preliminary.
Extension to C-MAPSS (bearing degradation, 21 sensors)~\cite{saxena2008damage}
and financial time series would strengthen the argument.

\textbf{No shared train/val windows:} the pipeline does not support
multi-source training (e.g., pre-train on seismic, fine-tune on 3W).
This is a deliberate design choice for cleaner evaluation, but limits
transfer learning scenarios.

\textbf{Hardware validation:} results are offline only.
The 8-channel Tiny TCN for 3W has a larger input layer and may require
more RAM than the univariate model; hardware validation on ESP32 is needed.

% ================================================================
\section{Conclusion}
\label{sec:conclusion}
% ================================================================

We evaluated a generic TinyML anomaly detection pipeline on two
structurally different domains---seismic (univariate, 40~Hz) and
Petrobras 3W (8-channel, 1~Hz)---using a unified dataset contract,
unified model suite, and unified metric set (AUC-PR, VUS-PR, PA-F1,
Event-F1, FP/h).

The adapter-based architecture enables this cross-domain evaluation
without any modification to training, MLOps, or export components.
Results show that [relative ranking stability / performance differences
to be discussed after filling TODO values].

The benchmark script, all adapters, profile files, and evaluation code
are publicly available as part of the companion pipeline repository.

Future work includes:
(i) adding C-MAPSS and a financial time series as additional benchmark
    domains;
(ii) multi-source training (seismic pre-training + 3W fine-tuning);
(iii) hardware validation of the multivariate Tiny TCN on ESP32.

% ================================================================
% REFERENCES
% ================================================================

\begin{thebibliography}{00}

\bibitem{pang2021deep}
G.~Pang et al.,
``Deep learning for anomaly detection: A review,''
\textit{ACM Comput. Surv.}, vol.~54, no.~2, pp. 1--38, 2021.

\bibitem{chandola2009anomaly}
V.~Chandola, A.~Banerjee, and V.~Kumar,
``Anomaly detection: A survey,''
\textit{ACM Comput. Surv.}, vol.~41, no.~3, pp. 1--58, 2009.

\bibitem{kim2022towards}
S.~Kim et al.,
``Towards a rigorous evaluation of time-series anomaly detection,''
in \textit{Proc. AAAI}, 2022.

\bibitem{vargas2019realistic}
R.~E.~V.~Vargas et al.,
``A realistic and public dataset with rare undesirable real events in oil wells,''
\textit{J. Petrol. Sci. Eng.}, vol.~181, p. 106223, 2019.

\bibitem{hundman2018detecting}
K.~Hundman et al.,
``Detecting spacecraft anomalies using LSTMs and nonparametric dynamic
thresholding,''
in \textit{Proc. KDD}, 2018.

\bibitem{su2019robust}
Y.~Su et al.,
``Robust anomaly detection for multivariate time series through stochastic
recurrent neural network,''
in \textit{Proc. KDD}, 2019.

\bibitem{shen2020timeseries}
L.~Shen et al.,
``Time series anomaly detection using temporal hierarchical one-class
network,''
in \textit{Proc. NeurIPS}, 2020.

\bibitem{bai2018empirical}
S.~Bai, J.~Z.~Kolter, and V.~Koltun,
``An empirical evaluation of generic convolutional and recurrent networks
for sequence modeling,''
arXiv:1803.01271, 2018.

\bibitem{saxena2008damage}
A.~Saxena et al.,
``Damage propagation modeling for aircraft engine run-to-failure simulation,''
in \textit{Proc. IEEE ICPHM}, 2008.

\bibitem{artigo1}
[AUTHOR],
``A generic TinyML pipeline for time-series anomaly detection,''
[VENUE], 2026. \textit{(companion paper 1)}

\bibitem{artigo2}
[AUTHOR],
``Operational quality gates for TinyML anomaly detection,''
[VENUE], 2026. \textit{(companion paper 2)}

\end{thebibliography}

\end{document}
